\subsection*{Important Caveat: Spectral Index Interpretation}

The predicted value $n_T \approx 4.90$ is extraordinarily blue. This is likely the single most problematic numerical claim in the series.

\subsubsection*{Why This Is Problematic}
Standard inflationary constraints require:
\begin{itemize}
    \item CMB bounds: $|n_T| \ll 1$ for scale-invariant spectra
    \item BBN: High-frequency GW overproduction constraints
    \item PTA: Upper bounds on stochastic GW backgrounds
    \item Energy density: $\Omega_{\text{GW}} \lesssim 10^{-6}$ at CMB scales
\end{itemize}

A persistent blue tensor spectrum with $n_T \approx 5$ would typically overproduce high-frequency gravitational radiation unless suppressed by:
\begin{itemize}
    \item A sharp cutoff at high $k$
    \item Transient generation epoch (not persistent)
    \item Narrow-band mechanism
    \item Modified transfer function
\end{itemize}

\subsubsection*{Proposed Reframing}
We therefore reframe $n_T$ as:
\begin{enumerate}
    \item A \textbf{transient effective exponent} within our truncated polynomial approximation
    \item \textbf{Not} an asymptotic inflationary tensor tilt
    \item Valid only within a limited range of scales $k_1 < k < k_2$
    \item Likely suppressed by a UV cutoff in a complete theory
\end{enumerate}

\subsubsection*{Revised Wording}
Replace:
\begin{quote}
"The spectral tilt $n_T$ is predicted to be $4 + 13\ln\phi \approx 4.90$."
\end{quote}
With:
\begin{quote}
"Within our truncation, the effective exponent governing scale dependence is $n_T^{\text{eff}} = 4 + 13\ln\phi \approx 4.90$. This is a \textit{transient} exponent, not an asymptotic inflationary tilt. A persistent blue spectrum would be observationally excluded; therefore we interpret this as a truncation artifact or as indicating a narrow-band generation mechanism."
\end{quote}
